\section{Introduction}

\subsection{Problem Statement and Motivation}

Artificial intelligence (AI) and, in particular, deep learning have enabled substantial advances in
medical signal analysis, supporting early detection, and continuous monitoring.

The focus of this Literature review is on newer deep learning models and other traditional machine learning for epileptic seizure detection and prediction using
non‑EEG wearable/autonomic biosignals (ECG/HRV, PPG, EDA, ACC). A key focus will be on
how the inclusion of time series data into the context of the model,
novel deep learning architectures, training strategies, and personalized fine tuning can improve performance
(sensitivity, false alarms per hour (FAR/h), latency) in clinical and ambulatory contexts.


Epilepsy affects roughly 7.6 per 1,000 people and about 30\% of
patients are drug resistant, motivating demand for timely, low‑burden
monitoring and early intervention \parencite{beghiEpidemiologyEpilepsy2019,
kwanEarlyIdentificationRefractory2000,chenNewEraPersonalised2020}.
While EEG remains the reference in controlled settings, EEG hardware
is obtrusive for pervasive use. Wearables offer the possibility of
continuous monitoring in ambulatory settings \parencite{chenNewEraPersonalised2020}. 
The key question is which model choices and architectures on wearable signals deliver clinically
viable performance outside the clinic.

\vspace{0.5em}
\noindent\textbf{Model‑centric challenges}
\begin{itemize}
  \item Operating point: achieving high sensitivity at acceptable
    FAR/h via personalized fine-tuning, calibration, thresholding, and alarm logic.
  \item Robustness and generalization: handling motion artifacts,
    physiological non‑stationarity, and domain shift (patient/device/site/activity).
  \item Interpretability and safety: explanation and fail‑safes suitable
    for regulated healthcare.
  \item Deployment context: clinic vs everyday use, sample‑ vs
    alarm‑based evaluation, latency/energy limits for on‑device
    inference, and the possibility for edge computing.
  \item Data limitations: scarce prospective, patient‑preserving,
    multimodal wearable datasets and the need for external/cross‑dataset
    validation, or real world testing.
\end{itemize}

\subsection{Evaluation Protocols (brief Overview)}
We distinguish the following evaluation elements:
\begin{enumerate}
  \item \textbf{Study designs:}
    \begin{itemize}
      \item Retrospective
      \item Pseudo‑prospective
      \item Prospective
    \end{itemize}
  \item \textbf{Evaluation metrics:}
    \begin{itemize}
      \item Sample‑based metrics
      \item Alarm‑based metrics
    \end{itemize}
  \item \textbf{Data splits:}
    \begin{itemize}
      \item Patient‑preserving splits
    \end{itemize}
\end{enumerate}
Alarm‑based sensitivity, false alarms per hour (FAR/h), and
latency better reflect real‑world usability than sample‑based accuracy and are key performance metrics.
Cross‑dataset transfer is essential to assess robustness and domain
shift \parencite{yangMultimodalAISystem2022,andradePerformanceSeizurePrediction2024}.

\subsection{Research Aim and Questions}
This literature review focuses on machine/deep learning models for seizure detection and prediction
using non-EEG wearable/autonomic biosignals, and determines the optimal strategies to achieve clinically
meaningful performance in ambulatory contexts.

\vspace{0.5em}
\subsubsection{Primary Research Question}
In non‑EEG wearable seizure detection and preictal prediction
(ECG/HRV, PPG, EDA, ACC), which modeling and evaluation choices
reported in 2015--2025 are associated with low alarm‑based false‑alarm
rates and high sensitivity with acceptable latency in ambulatory
settings?

\subsubsection{Sub‑questions}
\begin{enumerate}
  \item What ranges of sensitivity, FAR/h, and latency are reported
    under retrospective, pseudo‑prospective, and prospective designs,
    and how do alarm‑based results differ from sample‑based metrics?
  \item Across studies, how do traditional feature‑based models compare
    to 1D‑CNN/TCN/LSTM, and what evidence exists for attention/
    transformer variants on these signals, including any reported
    computational/latency considerations?
  \item How do temporal choices (window length, stride, context length)
    correlate with reported alarm‑based performance in real‑time use?
  \item What evidence is available on robustness to domain shift
    (patient/device/site/activity), including cross‑dataset/external
    validation and how do these affect alarm‑based metrics?
  \item What dataset gaps and class imbalance issues are reported, and
    what is the measured impact of augmentation or synthetic data on
    alarm‑based outcomes?
\end{enumerate}

\vspace{0.5em}
\noindent\textit{Contribution.} We synthesize a taxonomy of model
designs for wearable seizure detection/prediction, analyze how
evaluation choices affect reported performance, and derive practical
recommendations to achieve low FAR/h with high sensitivity in
ambulatory deployment.

\subsection{Methodological Approach and Structure}
The review follows structured IS/medical literature guidance (concept
matrix inspired by \cite{websterAnalyzingPastDetermine2002} and PRISMA
principles). The methodological approach uses a multi‑stage search and
transparent extraction/synthesis procedures as outlined below.

\begin{enumerate}
  \item \textbf{Search strategy (multi‑stage, 2015--2025):}
    \begin{itemize}
      \item Scoping: Google Scholar for broad coverage and identifying
        candidate papers.
      \item Formal queries: IEEE Xplore, PubMed, and Scopus using
        controlled keyword strings and deduplication.
      \item Citation chasing: backward and forward citation tracking on
        key papers to ensure completeness.
    \end{itemize}
  \item \textbf{Inclusion / Exclusion criteria:}
    \begin{itemize}
      \item Inclusion: human studies using wearable/non‑EEG autonomic
        signals for seizure detection or preictal prediction.
      \item Exclusion: EEG‑only studies, invasive recordings, or papers
        limited to purely algorithmic simulations without human data.
    \end{itemize}
\end{enumerate}

Extracted dimensions populate a concept matrix with the following
fields:
\begin{enumerate}
  \item Setting: clinical vs ambulatory.
  \item Modality: ECG/HRV, PPG, EDA, ACC.
  \item Task: detection vs prediction.
  \item Model class and fusion strategy.
  \item Study phase: retrospective, pseudo‑prospective, prospective.
  \item Evaluation granularity: sample‑ vs alarm‑based.
  \item Metrics: sensitivity, FAR/h, AUC, latency.
  \item Personalization.
  \item Interpretability and safety considerations.
  \item Data characteristics and dataset identity.
\end{enumerate}

\noindent\textbf{Quality flags:} patient‑preserving splits; external/
cross‑dataset validation; and FAR/h reporting. A PRISMA‑style flow
diagram will summarize screening and inclusion. Synthesis is narrative,
grouped by modality and task, with emphasis on model design levers and
deployment feasibility.

\section{Modeling landscape and design dimensions}

\subsection{Signals, tasks, and problem setup}
We target detection and preictal prediction from ECG/HRV, PPG, EDA,
and ACC. Windowing, labeling, and alarm logic define the operating
point; alarm‑based metrics (sensitivity, FAR/h, latency) are primary
for deployment.

\subsection{Representation learning and temporal modeling}
Classical models on HRV/EDA/ACC features (e.g., SVM, RF, kNN) achieve
solid detection and some prediction, often in retrospective designs
\parencite{sethFeasibilityCardiacbasedSeizure2023,masonHeartRateVariability2024}.
Deep architectures (1D‑CNN, LSTM, TCN, transformers) improve
robustness and phase discrimination, especially with attention
\parencite{yangMultimodalAISystem2022,pordoyEnhancedNonEEGMultimodal2025,
abtahiIdentificationRelevantECG2025,kalousiosECGbasedEpilepticSeizure2024}.

\subsection{Multimodal fusion}
Feature‑, decision‑, and representation‑level fusion can exploit
complementary ECG/EDA/ACC/PPG cues; attention‑based fusion often
yields gains under noise and non‑stationarity
\parencite{yangMultimodalAISystem2022,pordoyEnhancedNonEEGMultimodal2025}.

\subsection{Personalization, calibration, and drift}
Inter‑patient variability degrades population models; patient‑specific
fine‑tuning and post‑hoc calibration reduce FAR/h and improve
sensitivity in ambulatory, alarm‑based settings
\parencite{andradePerformanceSeizurePrediction2024}. Thresholding and
periodic recalibration mitigate drift.

\subsection{Robustness and domain shift}
Cross‑dataset evaluation (e.g., TUH→RPAH/EPILEPSIAE) exposes domain
shift and favors alarm‑ over sample‑based metrics
\parencite{yangMultimodalAISystem2022,andradePerformanceSeizurePrediction2024}.
Artifact handling, augmentation (time‑warping, jitter, scaling), and
OOD checks support generalization
\parencite{kalousiosECGbasedEpilepticSeizure2024,sethFeasibilityCardiacbasedSeizure2023}.

\subsection{Deployment considerations}
On‑device inference requires latency/energy awareness; compact
architectures and event‑driven processing are preferred. Interpretability
and alert justification are important for safety and clinical
acceptance.

\section{Technical background (clinical and sensing context)}

\subsection{Epileptic seizures (brief Overview)}
Seizures are classified by onset as focal (aware vs impaired
awareness; motor vs non‑motor, may evolve to bilateral tonic–clonic)
and generalized (motor: tonic–clonic, myoclonic, atonic; non‑motor:
absence) \parencite{fisherInstructionManualILAE2017}. Tonic–clonic
seizures commonly cause loss of consciousness and carry increased
injury/SUDEP risk \parencite{hesdorfferCombinedAnalysisRisk2011}.

\subsection{Autonomic markers and wearable sensing}
Sympathetic activation around seizures manifests as ictal tachycardia,
preictal HR rise with reduced HRV (lower RMSSD/HF, higher LF/HF),
increased EDA (tonic SCL, phasic SCRs), and limited ACC changes
useful for artifacts/context. PPG yields pulse‑rate variability
analogous to HRV from ECG \parencite{mironAutonomicBiosignalsSeizure2025,
masonHeartRateVariability2024}. Wearable heart‑rate sensing aligns
closely with chest‑strap ECG under motion (e.g., Apple Watch
Ultra/Polar H10 concordance \textasciitilde0.998), though dominant‑wrist
motion degrades accuracy; chest‑strap ECG provides robust RR intervals
during exercise \parencite{wolfHearttoWearAssessingAccuracy2024a}.

\section{Methodology}
\begin{enumerate}
  \item \textbf{Scope:} human studies in journals (2015--2025) on
    non‑EEG wearable/autonomic biosignals (ECG/HRV, PPG, EDA, ACC)
    for seizure detection or prediction.
  \item \textbf{Search:} Google Scholar (scoping), IEEE Xplore (all
    queries), PubMed and Scopus (targeted) using predefined query
    strings and deduplication.
  \item \textbf{Screening:} title/abstract screen, then full‑text
    eligibility; backward/forward citation chasing on key papers.
  \item \textbf{Extraction:} dataset type (wearable vs clinical), name,
    availability, task (detection vs prediction), evaluation design
    (retrospective/pseudo‑prospective/prospective), metrics
    (sensitivity, FAR/h, AUC, latency), personalization, cohort size/
    class balance.
  \item \textbf{Quality flags:} sample‑ vs alarm‑based evaluation;
    patient‑preserving splits; cross‑dataset/external validation.
  \item \textbf{Synthesis:} narrative grouping by modality and task with
    emphasis on modeling design levers and ambulatory feasibility.
\end{enumerate}

\section{Key Research Questions}
\begin{enumerate}
  \item Which wearable biosignals (ECG/HRV, EDA, ACC, PPG) and model
    classes (feature‑based, 1D‑CNN/LSTM/TCN, transformers, multimodal
    fusion) maintain performance under pseudo‑/prospective, alarm‑based
    evaluation?
  \item What performance gains arise from patient‑specific adaptation,
    and how can calibration/fine‑tuning be implemented under edge
    constraints?
  \item How can FAR/h be minimized via thresholding, calibration, and
    alarm logic while preserving high sensitivity in ambulatory use?
  \item How robust are models to domain shift (patient/device/site), and
    what is the impact of cross‑dataset/external validation and OOD
    checks?
  \item What dataset gaps exist for non‑EEG wearables, and how can
    augmentation/synthetic data mitigate imbalance without bias?
\end{enumerate}

\section{Overview of Literature Review}
This review is organized as follows.
\begin{enumerate}
  \item \textbf{Modeling landscape:} signals and tasks; representation
    learning; temporal modeling; fusion; personalization; robustness;
    deployment.
  \item \textbf{Clinical/sensing context (brief Overview):} seizure taxonomy; ANS
    markers; wearable sensing accuracy.
  \item \textbf{Methodology:} search, selection, extraction; PRISMA
    flowgraph.
  \item \textbf{Results:} concept matrix; comparative performance;
    modality/model analysis.
  \item \textbf{Discussion:} interpretation, trade‑offs, safety,
    deployment.
  \item \textbf{Limitations and future work.}
  \item \textbf{Conclusion.}
\end{enumerate}
